% --- IMPORTS ---
\documentclass[leqno]{article}
\usepackage{graphicx}
\usepackage{dsfont}
\usepackage{mathtools}
\usepackage{amssymb}
\usepackage{amsmath}
\usepackage{amsthm}
\usepackage{float}
\usepackage{ragged2e}
\usepackage{tensor}
\usepackage{fancyhdr}
\usepackage{empheq}
\usepackage[most]{tcolorbox}
\usepackage{enumitem}
\usepackage{dirtytalk}
\usepackage{marginnote}
\usepackage{microtype}
\usepackage[
  a4paper,
  left=0.8in,
  right=1in,
  top=1in,
  bottom=0.5in,
  headheight=14pt,
  headsep=0.4in
]{geometry}
\setlength{\parindent}{0cm}
% --- TikZ Packages ---
\usepackage{pgfplots}
\pgfplotsset{compat=1.18}
\usepgfplotslibrary{fillbetween, polar}
\usetikzlibrary{calc, positioning}
\usetikzlibrary{angles, quotes}
\usetikzlibrary{arrows.meta}
\usetikzlibrary{decorations.pathreplacing, decorations.markings}
\usetikzlibrary{patterns}
\usetikzlibrary{intersections}
% --- URL Packages ---
\usepackage[colorlinks=true,linkcolor=red,citecolor=magenta,urlcolor=black]{hyperref}%
\urlstyle{same}
% --- END OF IMPORTS ---

% --- CUSTOM COMMANDS ---
\RenewDocumentCommand{\marks}{o m}{%
  \IfValueTF{#1}%
    {\marginnote{\raggedright\textbf{[#2]}}[#1]}%
    {\marginnote{\raggedright\textbf{[#2]}}}%
}
\newcommand{\vspacerule}[2]{%
  \par\vspace*{#1}%
  \noindent\hrulefill\par
  \vspace*{#2}%
}
\definecolor{scarlet}{RGB}{205, 40, 75}
\newcommand{\questionitem}{%
  \item\phantomsection
  \addcontentsline{toc}{section}{Question \number\value{enumi}}%
}
\renewcommand{\contentsname}{Questions}
% --- END OF CUSTOM COMMANDS ---

% --- HEADER CONFIG ---
\pagestyle{fancy}
\fancyhead{}
\fancyfoot{}
\renewcommand{\headrulewidth}{0.9pt}
\lhead{\textit{Solutions} - Vector Product}
\chead{}
\rhead{\href{https://danieljsmith.org}{danieljsmith.org}}
% --- END OF HEADER CONFIG ---

\begin{document}
\vspace*{20pt}
\tableofcontents
\newpage
\begin{enumerate}[
  label=\textbf{\arabic*.},
  leftmargin=*,
  topsep=0pt,
  partopsep=0pt,
  parsep=0pt
]

% ---- Question 1 ----
\questionitem
% [Source: _QBT___Solns__Vector_Product, Q1]

\begin{enumerate}[label=\textbf{(\alph*)}]
  \item The vectors
  \[
  \mathbf{p}=\begin{pmatrix}1\\3\\2\end{pmatrix}
  \quad\text{and}\quad
  \mathbf{q}=\begin{pmatrix}2\\m\\-1\end{pmatrix}
  \]
  span a plane with normal vector
  \[
  \mathbf{n}=\begin{pmatrix}1\\-1\\1\end{pmatrix}
  \]
  Given that $\mathbf{q}\times\mathbf{p}$ is in the same direction as $\mathbf{n}$, find $m$ and the positive constant $\lambda$ such that
  \[
  \mathbf{q}\times\mathbf{p}=\lambda\mathbf{n}
  \]
  \marks[-20pt]{3}
  \item Using $\mathbf{q}\times\mathbf{p}$, determine the acute angle between $\mathbf{p}$ and $\mathbf{q}$. \marks{3} \\
\end{enumerate}

\vspace*{10pt}

\begin{tcolorbox}[colback=white, colframe=scarlet, title={\textbf{Solution}}, breakable, grow to left by=6mm, grow to right by=10mm]
\begin{enumerate}[label=\textbf{(\alph*)}]
  \item Taking
\[
\mathbf q=\begin{pmatrix}2\\m\\-1\end{pmatrix},
\qquad
\mathbf p=\begin{pmatrix}1\\3\\2\end{pmatrix}
\]
we find the cross product:
\begin{align*}
\mathbf q\times \mathbf p
&=
\begin{vmatrix}
\mathbf i & \mathbf j & \mathbf k\\
2 & m & -1\\
1 & 3 & 2
\end{vmatrix} \\
&=
\mathbf i(2m+3)-\mathbf j(4+1)+\mathbf k(6-m) \\
&=
\begin{pmatrix}
2m+3\\
-5\\
6-m
\end{pmatrix}
\end{align*}

Since $\mathbf q\times \mathbf p$ is in the same direction as
\[
\mathbf n=\begin{pmatrix}1\\-1\\1\end{pmatrix}
\]
we can write
\[
\mathbf q\times \mathbf p=\lambda \mathbf n
=
\begin{pmatrix}
\lambda\\
-\lambda\\
\lambda
\end{pmatrix}
\]

So
\[
\begin{pmatrix}
2m+3\\
-5\\
6-m
\end{pmatrix}
=
\begin{pmatrix}
\lambda\\
-\lambda\\
\lambda
\end{pmatrix}
\]

Equating components:
\[
-5=-\lambda \implies \lambda=5
\]

Then
\[
2m+3=5 \implies 2m=2 \implies m=1
\]

Checking the third component:
\[
6-m=6-1=5
\]
which is consistent.

Hence
\[
m=1,\qquad \lambda=5
\]

  \item Using $m=1$,
\[
\mathbf q=\begin{pmatrix}2\\1\\-1\end{pmatrix}
\]

Also
\[
\mathbf q\times \mathbf p=\begin{pmatrix}5\\-5\\5\end{pmatrix}
\]
so
\begin{align*}
|\mathbf q\times \mathbf p|
&=\sqrt{5^2+(-5)^2+5^2} \\
&=\sqrt{75} \\
&=5\sqrt3
\end{align*}

Now
\begin{align*}
|\mathbf p|&=\sqrt{1^2+3^2+2^2}=\sqrt{14} \\
|\mathbf q|&=\sqrt{2^2+1^2+(-1)^2}=\sqrt6
\end{align*}

Using
\[
|\mathbf q\times \mathbf p|=|\mathbf q|\,|\mathbf p|\sin\theta
\]
gives
\begin{align*}
\sin\theta
&=\frac{|\mathbf q\times \mathbf p|}{|\mathbf q|\,|\mathbf p|} \\
&=\frac{5\sqrt3}{\sqrt6\,\sqrt{14}} \\
&=\frac{5\sqrt3}{\sqrt{84}} \\
&=\frac{5}{2\sqrt7}
\end{align*}

So the acute angle is
\[
\theta=\sin^{-1}\!\left(\frac{5}{2\sqrt7}\right)\approx 70.9^\circ
\]

Hence the acute angle between $\mathbf p$ and $\mathbf q$ is
\[
\sin^{-1}\!\left(\frac{5}{2\sqrt7}\right)\approx 70.9^\circ
\]
\end{enumerate}
\end{tcolorbox}


\newpage

% ---- Question 2 ----
\questionitem
% [Source: _QBT___Solns__Vector_Product, Q2]

You are given the fixed points $A(1,0,1)$ and $C(4,0,2)$, and the variable point $B(t,2,5-t)$, where $t$ is a real parameter. \\

\begin{enumerate}[label=\textbf{(\alph*)}]
  \item Using only the geometrical properties of the vector product, explain why the statement ``$\overrightarrow{AB} \times \overrightarrow{AC} = 0$" is false for all values of $t$. \marks{2} \\
  \item
  \begin{enumerate}[label=(\roman*)]
    \item Use the vector product to find an expression, in terms of $t$, for the area of triangle $ABC$. \marks{4}
    \item Hence determine the value of $t$ for which the area of triangle $ABC$ is a minimum. \marks{2} \\
  \end{enumerate}
\end{enumerate}

\vspace*{10pt}

\begin{tcolorbox}[colback=white, colframe=scarlet, title={\textbf{Solution}}, breakable, grow to left by=6mm, grow to right by=10mm]
\begin{enumerate}[label=\textbf{(\alph*)}]
  \item 
  \[
  \overrightarrow{AB}=B-A=(t-1,\,2,\,4-t), \qquad \overrightarrow{AC}=C-A=(3,\,0,\,1)
  \]

  Geometrically, \(\overrightarrow{AB}\times\overrightarrow{AC}=\mathbf 0\) only if \(\overrightarrow{AB}\) and \(\overrightarrow{AC}\) are parallel, or if one of them is the zero vector.

  Here \(\overrightarrow{AC}=(3,0,1)\neq \mathbf 0\), and \(\overrightarrow{AB}\) can never be parallel to \(\overrightarrow{AC}\) because:

  \[
  \text{\(y\)-component of }\overrightarrow{AB}=2, \qquad \text{\(y\)-component of }\overrightarrow{AC}=0
  \]

  So no scalar multiple of \((3,0,1)\) can equal \((t-1,2,4-t)\) for any real \(t\).

  Therefore \(\overrightarrow{AB}\times\overrightarrow{AC}\neq \mathbf 0\) for all values of \(t\), so the statement is false for all \(t\).

  \item 
  \begin{enumerate}[label=(\roman*)]
    \item 
    First find the two sides of the triangle:

    \[
    \overrightarrow{AB}=(t-1,\,2,\,4-t), \qquad \overrightarrow{AC}=(3,\,0,\,1)
    \]

    Now calculate the vector product:

    \begin{align*}
    \overrightarrow{AB}\times\overrightarrow{AC}
    &=
    \begin{vmatrix}
    \mathbf i & \mathbf j & \mathbf k \\
    t-1 & 2 & 4-t \\
    3 & 0 & 1
    \end{vmatrix} \\
    &= \mathbf i\bigl(2\cdot 1-(4-t)\cdot 0\bigr)
    -\mathbf j\bigl((t-1)\cdot 1-(4-t)\cdot 3\bigr)
    +\mathbf k\bigl((t-1)\cdot 0-2\cdot 3\bigr) \\
    &= 2\mathbf i-\mathbf j\bigl(t-1-12+3t\bigr)-6\mathbf k \\
    &= 2\mathbf i-(4t-13)\mathbf j-6\mathbf k \\
    &= (2,\,13-4t,\,-6)
    \end{align*}

    The area of triangle \(ABC\) is

    \[
    \frac12\left|\overrightarrow{AB}\times\overrightarrow{AC}\right|
    \]

    So

    \begin{align*}
    \text{Area}
    &= \frac12\sqrt{2^2+(13-4t)^2+(-6)^2} \\
    &= \frac12\sqrt{4+(13-4t)^2+36} \\
    &= \frac12\sqrt{(13-4t)^2+40}
    \end{align*}

    Hence the area of triangle \(ABC\) is

    \[
    \frac12\sqrt{(13-4t)^2+40}
    \]

    \item 
    To minimise the area, it is enough to minimise

    \[
    (13-4t)^2+40
    \]

    Since \((13-4t)^2\geq 0\), its smallest value is \(0\), which occurs when

    \[
    13-4t=0
    \]

    Hence

    \[
    t=\frac{13}{4}
    \]

    The minimum area is then

    \[
    \frac12\sqrt{40}=\sqrt{10}
    \]

    Therefore the area is minimised when \(t=\dfrac{13}{4}\), and the minimum area is \(\sqrt{10}\).
  \end{enumerate}
\end{enumerate}
\end{tcolorbox}


\newpage

% ---- Question 3 ----
\questionitem
% [Source: _QBT___Solns__Vector_Product, Q3]

With respect to a fixed origin $O$, the points $A$ and $B$ have coordinates $(1, 1, 0)$ and $(p, 3, 2)$ respectively, where $p$ is a constant. \\

For each of the following, determine the possible values of $p$ for which \\

\begin{enumerate}[label=\textbf{(\alph*)}]
  \item $\overrightarrow{AB}$ makes an angle of $60^\circ$ with the positive $y$-axis. \marks{3} \\
  \item $\overrightarrow{OA} \times \overrightarrow{AB}$ is parallel to
  \[
  \begin{pmatrix}2\\-2\\1\end{pmatrix}
  \]
  \marks[-20pt]{3}
  \item the area of triangle $OAB$ is $\tfrac{3}{2}$. \marks{3} \\
\end{enumerate}

\vspace*{10pt}

\begin{tcolorbox}[colback=white, colframe=scarlet, title={\textbf{Solution}}, breakable, grow to left by=6mm, grow to right by=10mm]
\begin{enumerate}[label=\textbf{(\alph*)}]
  \item
  First find the vector
  \[
  \overrightarrow{AB}
  =\begin{pmatrix}p-1\\3-1\\2-0\end{pmatrix}
  =\begin{pmatrix}p-1\\2\\2\end{pmatrix}
  \]

  The positive \(y\)-axis has direction vector
  \[
  \mathbf{j}=\begin{pmatrix}0\\1\\0\end{pmatrix}
  \]

  Using the scalar product formula,
  \[
  \overrightarrow{AB}\cdot \mathbf{j}
  =\left|\overrightarrow{AB}\right|\left|\mathbf{j}\right|\cos 60^\circ
  \]

  So
  \begin{align*}
  2&=\sqrt{(p-1)^2+2^2+2^2}\times 1\times \frac12 \\
  2&=\frac12\sqrt{(p-1)^2+8} \\
  4&=\sqrt{(p-1)^2+8} \\
  16&=(p-1)^2+8 \\
  (p-1)^2&=8
  \end{align*}

  Hence
  \[
  p-1=\pm 2\sqrt{2}
  \]

  Therefore the possible values of \(p\) are
  \[
  p=1\pm 2\sqrt{2}
  \]

  \item
  \[
  \overrightarrow{OA}=\begin{pmatrix}1\\1\\0\end{pmatrix},
  \qquad
  \overrightarrow{AB}=\begin{pmatrix}p-1\\2\\2\end{pmatrix}
  \]

  Now calculate the cross product:
  \begin{align*}
  \overrightarrow{OA}\times \overrightarrow{AB}
  &=
  \begin{vmatrix}
  \mathbf{i}&\mathbf{j}&\mathbf{k}\\
  1&1&0\\
  p-1&2&2
  \end{vmatrix} \\
  &=\mathbf{i}(1\cdot 2-0\cdot 2)-\mathbf{j}(1\cdot 2-0\cdot (p-1))+\mathbf{k}(1\cdot 2-1\cdot (p-1)) \\
  &=\begin{pmatrix}2\\-2\\3-p\end{pmatrix}
  \end{align*}

  For this to be parallel to \(\begin{pmatrix}2\\-2\\1\end{pmatrix}\), we must have
  \[
  \begin{pmatrix}2\\-2\\3-p\end{pmatrix}
  =\lambda \begin{pmatrix}2\\-2\\1\end{pmatrix}
  \]

  Comparing first components,
  \[
  2=2\lambda \implies \lambda=1
  \]

  Then comparing third components,
  \[
  3-p=1
  \]

  So
  \[
  p=2
  \]

  \item
  For triangle \(OAB\), the area is
  \[
  \text{Area}=\frac12\left|\overrightarrow{OA}\times \overrightarrow{OB}\right|
  \]

  Since
  \[
  \overrightarrow{OB}=\overrightarrow{OA}+\overrightarrow{AB}
  \]

  we have
  \begin{align*}
  \overrightarrow{OA}\times \overrightarrow{OB}
  &=\overrightarrow{OA}\times (\overrightarrow{OA}+\overrightarrow{AB}) \\
  &=\overrightarrow{OA}\times \overrightarrow{OA}+\overrightarrow{OA}\times \overrightarrow{AB} \\
  &=\mathbf{0}+\overrightarrow{OA}\times \overrightarrow{AB} \\
  &=\begin{pmatrix}2\\-2\\3-p\end{pmatrix}
  \end{align*}

  Therefore
  \begin{align*}
  \frac12\sqrt{2^2+(-2)^2+(3-p)^2}&=\frac32 \\
  \frac12\sqrt{8+(3-p)^2}&=\frac32 \\
  \sqrt{8+(3-p)^2}&=3 \\
  8+(3-p)^2&=9 \\
  (3-p)^2&=1
  \end{align*}

  So
  \[
  3-p=\pm 1
  \]

  Hence the possible values of \(p\) are
  \[
  p=2 \text{ or } p=4
  \]
\end{enumerate}
\end{tcolorbox}


\newpage

% ---- Question 4 ----
\questionitem
% [Source: _QBT___Solns__Vector_Product, Q4]

Points $P$, $Q$ and $R$ have position vectors $\mathbf{p}$, $\mathbf{q}$ and $\mathbf{r}$ respectively, relative to origin $O$. \\

It is given that $(\mathbf{q}-\mathbf{p}) \times (\mathbf{r}-\mathbf{q}) = 2\mathbf{p}$ and that $\lvert \mathbf{p} \rvert = 5$. \\

\begin{enumerate}[label=\textbf{(\alph*)}]
  \item Determine each of the following as either a single vector or a scalar quantity.
  \begin{enumerate}[label=(\roman*)]
    \item $(\mathbf{r}-\mathbf{q}) \times (\mathbf{q}-\mathbf{p})$ \marks{1}
    \item $\mathbf{p} \times \big((\mathbf{q}-\mathbf{p}) \times (\mathbf{r}-\mathbf{q})\big)$ \marks{2}
    \item $\mathbf{p} \cdot \big((\mathbf{q}-\mathbf{p}) \times (\mathbf{r}-\mathbf{q})\big)$ \marks{2} \\
  \end{enumerate}
  \item Describe a geometrical relationship between the points $O$, $P$, $Q$ and $R$ which can be deduced from
  \begin{enumerate}[label=(\roman*)]
    \item the statement $(\mathbf{q}-\mathbf{p}) \times (\mathbf{r}-\mathbf{q}) = 2\mathbf{p}$ \marks{1}
    \item the result of part (a)(iii) \marks{1} \\
  \end{enumerate}
\end{enumerate}

\vspace*{10pt}

\begin{tcolorbox}[colback=white, colframe=scarlet, title={\textbf{Solution}}, breakable, grow to left by=6mm, grow to right by=10mm]
\begin{enumerate}[label=\textbf{(\alph*)}]
  \item
  \begin{enumerate}[label=(\roman*)]
    \item Using the fact that $\mathbf a\times\mathbf b=-\mathbf b\times\mathbf a$,
    \[
    (\mathbf r-\mathbf q)\times(\mathbf q-\mathbf p)
    =-\big((\mathbf q-\mathbf p)\times(\mathbf r-\mathbf q)\big)
    =-2\mathbf p
    \]
    So the vector is $-2\mathbf p$.

    \item Substitute the given result directly:
    \[
    \mathbf p\times\big((\mathbf q-\mathbf p)\times(\mathbf r-\mathbf q)\big)
    =\mathbf p\times 2\mathbf p
    =2(\mathbf p\times\mathbf p)
    =\mathbf 0
    \]
    So the vector is $\mathbf 0$.

    \item Again using the given result,
    \begin{align*}
    \mathbf p\cdot\big((\mathbf q-\mathbf p)\times(\mathbf r-\mathbf q)\big)
    &=\mathbf p\cdot 2\mathbf p\\
    &=2(\mathbf p\cdot\mathbf p)\\
    &=2\lvert\mathbf p\rvert^2\\
    &=2(5^2)\\
    &=50
    \end{align*}
    So the scalar is $50$.
  \end{enumerate}

  \item
  \begin{enumerate}[label=(\roman*)]
    \item The vector $\mathbf q-\mathbf p$ is in the direction of $PQ$, and $\mathbf r-\mathbf q$ is in the direction of $QR$.

    Their cross product is perpendicular to both of these vectors. Since
    \[
    (\mathbf q-\mathbf p)\times(\mathbf r-\mathbf q)=2\mathbf p
    \]
    the vector $\mathbf p$ is parallel to this perpendicular direction.

    But $\mathbf p$ is the position vector of $P$, so it is the direction of $OP$. Therefore $OP$ is perpendicular to both $PQ$ and $QR$, so $OP$ is normal to the plane $PQR$.

    Hence $OP\perp$ plane $PQR$.

    \item From part (a)(iii),
    \[
    \mathbf p\cdot\big((\mathbf q-\mathbf p)\times(\mathbf r-\mathbf q)\big)=50\neq 0
    \]
    A non-zero scalar triple product means the three vectors involved are not coplanar.

    Therefore the points $O$, $P$, $Q$ and $R$ are not coplanar.
  \end{enumerate}
\end{enumerate}
\end{tcolorbox}


\newpage

% ---- Question 5 ----
\questionitem
% [Source: _QBT___Solns__Vector_Product, Q5]

\begin{enumerate}[label=\textbf{(\alph*)}]
  \item Let
  \[
  \mathbf{p}=\begin{pmatrix}1\\1\\2\end{pmatrix}
  \quad\text{and}\quad
  \mathbf{q}=\begin{pmatrix}x\\0\\1\end{pmatrix}
  \]
  The area of the parallelogram formed by $\mathbf{p}$ and $\mathbf{q}$ is $\sqrt{2}$.\\
  It is also given that $\mathbf{p}\times\mathbf{q}$ is perpendicular to
  \[
  \begin{pmatrix}1\\3\\m\end{pmatrix}
  \]
  Find $x$ and $m$. \marks{4} \\
  \item Hence, using $|\mathbf{p}\times\mathbf{q}|$, determine the acute angle between $\mathbf{p}$ and $\mathbf{q}$. \marks{3} \\
\end{enumerate}

\vspace*{10pt}

\begin{tcolorbox}[colback=white, colframe=scarlet, title={\textbf{Solution}}, breakable, grow to left by=6mm, grow to right by=10mm]
\begin{enumerate}[label=\textbf{(\alph*)}]
  \item For two vectors, the area of the parallelogram they form is
\[
|\mathbf p\times \mathbf q|
\]

First find \(\mathbf p\times \mathbf q\):
\begin{align*}
\mathbf p\times \mathbf q
&=
\begin{vmatrix}
\mathbf i & \mathbf j & \mathbf k\\
1&1&2\\
x&0&1
\end{vmatrix} \\
&= \mathbf i(1\cdot 1-2\cdot 0)-\mathbf j(1\cdot 1-2x)+\mathbf k(1\cdot 0-1\cdot x) \\
&= \mathbf i-\mathbf j(1-2x)-x\mathbf k \\
&= \begin{pmatrix}1\\2x-1\\-x\end{pmatrix}
\end{align*}

Since the area is \(\sqrt2\),
\begin{align*}
|\mathbf p\times \mathbf q|&=\sqrt2 \\
\therefore \ |\mathbf p\times \mathbf q|^2&=2
\end{align*}

So
\begin{align*}
1^2+(2x-1)^2+(-x)^2&=2 \\
1+(4x^2-4x+1)+x^2&=2 \\
5x^2-4x&=0 \\
x(5x-4)&=0
\end{align*}

Hence
\[
x=0 \quad \text{or} \quad x=\frac45
\]

Now use the condition that \(\mathbf p\times \mathbf q\) is perpendicular to \(\begin{pmatrix}1\\3\\m\end{pmatrix}\).

Perpendicular vectors have zero dot product, so
\begin{align*}
\begin{pmatrix}1\\2x-1\\-x\end{pmatrix}\cdot \begin{pmatrix}1\\3\\m\end{pmatrix}&=0 \\
1+3(2x-1)+(-x)m&=0 \\
1+6x-3-mx&=0 \\
x(6-m)&=2
\end{align*}

This shows \(x\neq 0\), so the value \(x=0\) must be rejected.

Therefore
\[
x=\frac45
\]

Substitute into \(x(6-m)=2\):
\begin{align*}
\frac45(6-m)&=2 \\
6-m&=\frac{5}{2} \\
m&=6-\frac52=\frac72
\end{align*}

Therefore,
\[
x=\frac45,\qquad m=\frac72
\]

  \item Using
\[
|\mathbf p\times \mathbf q|=|\mathbf p|\,|\mathbf q|\sin\theta
\]
where \(\theta\) is the angle between \(\mathbf p\) and \(\mathbf q\).

From part (a),
\[
|\mathbf p\times \mathbf q|=\sqrt2
\]

Also,
\begin{align*}
|\mathbf p|
&=\sqrt{1^2+1^2+2^2}
=\sqrt6
\end{align*}

and with \(x=\frac45\),
\begin{align*}
|\mathbf q|
&=\sqrt{\left(\frac45\right)^2+0^2+1^2} \\
&=\sqrt{\frac{16}{25}+1}
=\sqrt{\frac{41}{25}}
=\frac{\sqrt{41}}{5}
\end{align*}

So
\begin{align*}
\sqrt2
&=\sqrt6\cdot \frac{\sqrt{41}}{5}\sin\theta \\
\sin\theta
&=\frac{\sqrt2}{\sqrt6(\sqrt{41}/5)} \\
&=\frac{5\sqrt2}{\sqrt{246}} \\
&=\frac{5}{\sqrt{123}}
\end{align*}

Hence the acute angle is
\[
\theta=\sin^{-1}\left(\frac{5}{\sqrt{123}}\right)\approx 26.8^\circ
\]

Therefore, the acute angle between \(\mathbf p\) and \(\mathbf q\) is
\[
\sin^{-1}\left(\frac{5}{\sqrt{123}}\right)\approx 26.8^\circ
\]
\end{enumerate}
\end{tcolorbox}


\newpage

% ---- Question 6 ----
\questionitem
% [Source: _QBT___Solns__Vector_Product, Q6]

The point $P$ has coordinates $P(2,-1,1)$. The vector $\overrightarrow{PQ}=(1,2,0)$ and the midpoint of $PR$ is $M(1,0,2)$. \\

\begin{enumerate}[label=\textbf{(\alph*)}]
  \item Use a vector product to show that the area of triangle $PQR$ is $\sqrt{14}$. \marks{4} \\
  \item Find a vector equation of the plane containing $P$, $Q$ and $R$ in the form $\mathbf{r} \cdot \mathbf{n} = k$. \marks{1} \\
  \item Hence find the exact distance of the plane from the origin. \marks{1} \\
\end{enumerate}

\vspace*{10pt}

\begin{tcolorbox}[colback=white, colframe=scarlet, title={\textbf{Solution}}, breakable, grow to left by=6mm, grow to right by=10mm]
\begin{enumerate}[label=\textbf{(\alph*)}]
  \item First find the coordinates of $Q$ and $R$.

  Since $\overrightarrow{PQ}=(1,2,0)$,
  \[
  Q=P+\overrightarrow{PQ}=(2,-1,1)+(1,2,0)=(3,1,1)
  \]

  Since $M(1,0,2)$ is the midpoint of $PR$,
  \[
  M=\frac{1}{2}(P+R)
  \]
  so
  \[
  R=2M-P=(2,0,4)-(2,-1,1)=(0,1,3)
  \]

  Now
  \[
  \overrightarrow{PR}=R-P=(0,1,3)-(2,-1,1)=(-2,2,2)
  \]

  Using the vector product,
  \[
  \overrightarrow{PQ}\times\overrightarrow{PR}
  =
  \begin{pmatrix}
  1\\
  2\\
  0
  \end{pmatrix}
  \times
  \begin{pmatrix}
  -2\\
  2\\
  2
  \end{pmatrix}
  =
  \begin{pmatrix}
  2\cdot 2-0\cdot 2\\
  0\cdot(-2)-1\cdot 2\\
  1\cdot 2-2\cdot(-2)
  \end{pmatrix}
  =
  \begin{pmatrix}
  4\\
  -2\\
  6
  \end{pmatrix}
  \]

  Its magnitude is
  \[
  \left|\overrightarrow{PQ}\times\overrightarrow{PR}\right|
  =\sqrt{4^2+(-2)^2+6^2}
  =\sqrt{16+4+36}
  =\sqrt{56}
  =2\sqrt{14}
  \]

  The area of triangle $PQR$ is half the area of the parallelogram:
  \[
  \text{Area of } \triangle PQR=\frac{1}{2}\left|\overrightarrow{PQ}\times\overrightarrow{PR}\right|
  =\frac{1}{2}(2\sqrt{14})
  =\sqrt{14}
  \]

  Hence the area of triangle $PQR$ is $\sqrt{14}$.

  \item A normal vector to the plane is
  \[
  \mathbf n=
  \begin{pmatrix}
  4\\
  -2\\
  6
  \end{pmatrix}
  \]
  which can be simplified to
  \[
  \mathbf n=
  \begin{pmatrix}
  2\\
  -1\\
  3
  \end{pmatrix}
  \]

  Using the point $P(2,-1,1)$,
  \[
  k=\mathbf p\cdot \mathbf n
  =
  \begin{pmatrix}
  2\\
  -1\\
  1
  \end{pmatrix}
  \cdot
  \begin{pmatrix}
  2\\
  -1\\
  3
  \end{pmatrix}
  =4+1+3=8
  \]

  So the plane is
  \[
  \mathbf r\cdot
  \begin{pmatrix}
  2\\
  -1\\
  3
  \end{pmatrix}
  =8
  \]

  \item For a plane $\mathbf r\cdot\mathbf n=d$, the distance from a point $\mathbf p$ is
  \[
  \mathrm{dist}=\frac{\left|\mathbf p\cdot\mathbf n-d\right|}{|\mathbf n|}
  \]

  Here the point is the origin, so $\mathbf p=\mathbf 0$, and the plane is
  \[
  \mathbf r\cdot
  \begin{pmatrix}
  2\\
  -1\\
  3
  \end{pmatrix}
  =8
  \]

  Therefore
  \[
  \mathrm{dist}
  =
  \frac{|0-8|}{\sqrt{2^2+(-1)^2+3^2}}
  =
  \frac{8}{\sqrt{14}}
  =
  \frac{4\sqrt{14}}{7}
  \]

  Hence the exact distance of the plane from the origin is $\dfrac{4\sqrt{14}}{7}$.
\end{enumerate}
\end{tcolorbox}


\newpage

% ---- Question 7 ----
\questionitem
% [Source: _QBT___Solns__Vector_Product, Q7]

The points $P$, $Q$ and $R$ have position vectors $\mathbf{p}=2\mathbf{i}-2\mathbf{k}$, $\mathbf{q}=3\mathbf{j}+3\mathbf{k}$ and $\mathbf{r}=t\mathbf{i}-\mathbf{j}+\mathbf{k}$ respectively, relative to the origin $O$. \\

Determine the value(s) of $t$ in each of the following cases. \\

\begin{enumerate}[label=\textbf{(\alph*)}]
  \item The line $OR$ is parallel to a normal to the plane $OPQ$. \marks{2} \\
  \item The volume of tetrahedron $OPQR$ is $9$. \marks{4} \\
\end{enumerate}

\vspace*{10pt}

\begin{tcolorbox}[colback=white, colframe=scarlet, title={\textbf{Solution}}, breakable, grow to left by=6mm, grow to right by=10mm]
\begin{enumerate}[label=\textbf{(\alph*)}]
  \item The plane $OPQ$ contains the vectors
\[
\mathbf p=\begin{pmatrix}2\\0\\-2\end{pmatrix},
\qquad
\mathbf q=\begin{pmatrix}0\\3\\3\end{pmatrix}
\]

A normal to the plane is $\mathbf p\times \mathbf q$:
\begin{align*}
\mathbf p\times \mathbf q
&=
\begin{vmatrix}
\mathbf i & \mathbf j & \mathbf k\\
2&0&-2\\
0&3&3
\end{vmatrix} \\
&= \mathbf i(0\cdot 3-(-2)\cdot 3)-\mathbf j(2\cdot 3-(-2)\cdot 0)+\mathbf k(2\cdot 3-0\cdot 0) \\
&= 6\mathbf i-6\mathbf j+6\mathbf k \\
&= 6(\mathbf i-\mathbf j+\mathbf k)
\end{align*}

The line $OR$ has direction vector
\[
\mathbf r=t\mathbf i-\mathbf j+\mathbf k
\]

If $OR$ is parallel to a normal to the plane, then $\mathbf r$ must be parallel to $\mathbf i-\mathbf j+\mathbf k$.

So
\[
(t,-1,1)\parallel (1,-1,1)
\]

Comparing components gives $t=1$.

Hence the value of $t$ is $1$.

  \item Using the vector formula for the volume of a tetrahedron,
\[
\text{Volume of }OPQR=\frac16\left|\mathbf p\cdot(\mathbf q\times \mathbf r)\right|
\]

Equivalently,
\[
\text{Volume}=\frac16\left|(\mathbf p\times \mathbf q)\cdot \mathbf r\right|
\]

From part (a),
\[
\mathbf p\times \mathbf q=6(\mathbf i-\mathbf j+\mathbf k)
\]

Also,
\[
\mathbf r=t\mathbf i-\mathbf j+\mathbf k
\]

So the scalar triple product is
\begin{align*}
(\mathbf p\times \mathbf q)\cdot \mathbf r
&=6(\mathbf i-\mathbf j+\mathbf k)\cdot (t\mathbf i-\mathbf j+\mathbf k) \\
&=6(t+1+1) \\
&=6(t+2)
\end{align*}

Therefore
\begin{align*}
\text{Volume}
&=\frac16|6(t+2)| \\
&=|t+2|
\end{align*}

Given that the volume is $9$,
\[
|t+2|=9
\]

So
\begin{align*}
t+2&=9 \quad \text{or} \quad t+2=-9 \\
t&=7 \quad \text{or} \quad t=-11
\end{align*}

Hence the values of $t$ are $7$ and $-11$.
\end{enumerate}
\end{tcolorbox}


\newpage

% ---- Question 8 ----
\questionitem
% [Source: _QBT___Solns__Vector_Product, Q8]

The points $A$, $B$ and $C$ lie on the plane $\Pi$. \\

The position vectors of $A$ and $B$ are $\mathbf{a}=2\mathbf{i}-\mathbf{j}+3\mathbf{k}$ and $\mathbf{b}=5\mathbf{i}+\mathbf{j}$ respectively. \\

The centroid $G$ of triangle $ABC$ has position vector $4\mathbf{i}+\mathbf{j}+2\mathbf{k}$. \\

\begin{enumerate}[label=\textbf{(\alph*)}]
  \item Find $\overrightarrow{AB} \times \overrightarrow{AC}$. \marks{4} \\
  \item Find an equation for $\Pi$ in the form $\mathbf{r} \cdot \mathbf{n} = p$. \marks{2} \\
\end{enumerate}
The point $D$ has position vector $\mathbf{i}+5\mathbf{j}+4\mathbf{k}$. \\
\begin{enumerate}[label=\textbf{(\alph*)}, resume]
  \item Determine the volume of the tetrahedron $ABCD$. \marks{4} \\
\end{enumerate}

\vspace*{10pt}

\begin{tcolorbox}[colback=white, colframe=scarlet, title={\textbf{Solution}}, breakable, grow to left by=6mm, grow to right by=10mm]
\begin{enumerate}[label=\textbf{(\alph*)}]
  \item The centroid of triangle \(ABC\) has position vector
\[
\mathbf{g}=\frac{\mathbf{a}+\mathbf{b}+\mathbf{c}}{3}
\]
So
\[
\mathbf{c}=3\mathbf{g}-\mathbf{a}-\mathbf{b}
\]

Using
\[
\mathbf{a}=\begin{pmatrix}2\\-1\\3\end{pmatrix},
\quad
\mathbf{b}=\begin{pmatrix}5\\1\\0\end{pmatrix},
\quad
\mathbf{g}=\begin{pmatrix}4\\1\\2\end{pmatrix}
\]
gives
\begin{align*}
\mathbf{c}
&=3\begin{pmatrix}4\\1\\2\end{pmatrix}
-\begin{pmatrix}2\\-1\\3\end{pmatrix}
-\begin{pmatrix}5\\1\\0\end{pmatrix} \\
&=\begin{pmatrix}12\\3\\6\end{pmatrix}
-\begin{pmatrix}2\\-1\\3\end{pmatrix}
-\begin{pmatrix}5\\1\\0\end{pmatrix} \\
&=\begin{pmatrix}5\\3\\3\end{pmatrix}
\end{align*}

Now
\begin{align*}
\overrightarrow{AB}
&=\mathbf{b}-\mathbf{a}
=\begin{pmatrix}5\\1\\0\end{pmatrix}-\begin{pmatrix}2\\-1\\3\end{pmatrix}
=\begin{pmatrix}3\\2\\-3\end{pmatrix} \\
\overrightarrow{AC}
&=\mathbf{c}-\mathbf{a}
=\begin{pmatrix}5\\3\\3\end{pmatrix}-\begin{pmatrix}2\\-1\\3\end{pmatrix}
=\begin{pmatrix}3\\4\\0\end{pmatrix}
\end{align*}

Hence
\begin{align*}
\overrightarrow{AB}\times\overrightarrow{AC}
&=
\begin{vmatrix}
\mathbf{i} & \mathbf{j} & \mathbf{k} \\
3 & 2 & -3 \\
3 & 4 & 0
\end{vmatrix} \\
&=\mathbf{i}(2\cdot 0-(-3)\cdot 4)-\mathbf{j}(3\cdot 0-(-3)\cdot 3)+\mathbf{k}(3\cdot 4-2\cdot 3) \\
&=12\mathbf{i}-9\mathbf{j}+6\mathbf{k}
\end{align*}

So
\[
\overrightarrow{AB}\times\overrightarrow{AC}=12\mathbf{i}-9\mathbf{j}+6\mathbf{k}
\]

  \item A normal vector to the plane is \(\overrightarrow{AB}\times\overrightarrow{AC}\), so we may take
\[
\mathbf{n}=12\mathbf{i}-9\mathbf{j}+6\mathbf{k}=3(4\mathbf{i}-3\mathbf{j}+2\mathbf{k})
\]
A simpler normal vector is
\[
\mathbf{n}=4\mathbf{i}-3\mathbf{j}+2\mathbf{k}
\]

The plane has equation \(\mathbf{r}\cdot\mathbf{n}=p\). Since \(A\) lies on the plane,
\begin{align*}
p&=\mathbf{a}\cdot\mathbf{n} \\
&=(2,-1,3)\cdot(4,-3,2) \\
&=2(4)+(-1)(-3)+3(2) \\
&=8+3+6=17
\end{align*}

Therefore the equation of \(\Pi\) is
\[
\mathbf{r}\cdot(4\mathbf{i}-3\mathbf{j}+2\mathbf{k})=17
\]

  \item First find
\[
\mathbf{d}=\begin{pmatrix}1\\5\\4\end{pmatrix}
\]
so
\begin{align*}
\overrightarrow{AD}
&=\mathbf{d}-\mathbf{a}
=\begin{pmatrix}1\\5\\4\end{pmatrix}-\begin{pmatrix}2\\-1\\3\end{pmatrix}
=\begin{pmatrix}-1\\6\\1\end{pmatrix}
\end{align*}

The volume of a tetrahedron with three edges \(\overrightarrow{AB}\), \(\overrightarrow{AC}\), \(\overrightarrow{AD}\) meeting at \(A\) is
\[
V=\frac{1}{6}\left|(\overrightarrow{AB}\times\overrightarrow{AC})\cdot\overrightarrow{AD}\right|
\]

Using the result from part (a),
\begin{align*}
V
&=\frac{1}{6}\left|(12,-9,6)\cdot(-1,6,1)\right| \\
&=\frac{1}{6}\left|12(-1)+(-9)(6)+6(1)\right| \\
&=\frac{1}{6}\left|-12-54+6\right| \\
&=\frac{1}{6}\left|-60\right| \\
&=10
\end{align*}

So the volume of tetrahedron \(ABCD\) is
\(10\)
\end{enumerate}
\end{tcolorbox}


\newpage

% ---- Question 9 ----
\questionitem
% [Source: _QBT___Solns__Vector_Product, Q9]

With respect to a fixed origin $O$ the point $A$ has coordinates $(9,10,-1)$ and the line $\ell$ is the intersection of the planes $x-z=-2$ and $x+2y=5$. \\

The points $B$ and $C$ lie on $\ell$ such that $AB=AC=15$. \\

Given that $A$ does not lie on $\ell$ and that the $z$ coordinate of $B$ is negative \\

\begin{enumerate}[label=\textbf{(\alph*)}]
  \item determine the coordinates of $B$ and the coordinates of $C$. \marks{4} \\
  \item Hence determine a Cartesian equation of the plane containing the points $A$, $B$ and $C$. \marks{3} \\
\end{enumerate}
The point $D$ has coordinates $(8,4,\alpha)$, where $\alpha$ is a constant. \\

Given that the volume of the tetrahedron $ABCD$ is $156$ \\

\begin{enumerate}[label=\textbf{(\alph*)}, resume]
  \item determine the possible values of $\alpha$. \marks{4} \\
\end{enumerate}
Given that $\alpha>0$ \\

\begin{enumerate}[label=\textbf{(\alph*)}, resume]
  \item determine the shortest distance between the line $\ell$ and the line passing through the points $A$ and $D$, giving your answer to $2$ significant figures. \marks{4} \\
\end{enumerate}

\vspace*{10pt}

\begin{tcolorbox}[colback=white, colframe=scarlet, title={\textbf{Solution}}, breakable, grow to left by=6mm, grow to right by=10mm]
\begin{enumerate}[label=\textbf{(\alph*)}]
\item Since $l$ is the intersection of
\[
x-z=-2
\qquad\text{and}\qquad
x+2y=5
\]
let $z=t$. Then
\[
x=t-2
\]
and substituting into $x+2y=5$ gives
\[
t-2+2y=5
\]
so
\[
2y=7-t
\qquad\Rightarrow\qquad
y=\frac{7-t}{2}
\]
Hence a general point $P$ on $l$ is
\[
P\left(t-2,\frac{7-t}{2},t\right)
\]

Now $AB=AC=15$, so points $B$ and $C$ on $l$ satisfy $AP=15$.

Using $A=(9,10,-1)$,
\begin{align*}
AP^2
&=(t-2-9)^2+\left(\frac{7-t}{2}-10\right)^2+(t-(-1))^2\\
&=(t-11)^2+\left(\frac{-t-13}{2}\right)^2+(t+1)^2
\end{align*}
Since $AP=15$,
\[
(t-11)^2+\left(\frac{t+13}{2}\right)^2+(t+1)^2=225
\]
Multiply by $4$:
\begin{align*}
4(t-11)^2+(t+13)^2+4(t+1)^2&=900\\
4(t^2-22t+121)+(t^2+26t+169)+4(t^2+2t+1)&=900\\
9t^2-54t+657&=900\\
9t^2-54t-243&=0\\
t^2-6t-27&=0\\
(t-9)(t+3)&=0
\end{align*}
So
\[
t=9 \quad \text{or} \quad t=-3
\]

If $t=-3$,
\[
P=\left(-5,5,-3\right)
\]
If $t=9$,
\[
P=\left(7,-1,9\right)
\]

Given that the $z$-coordinate of $B$ is negative, we take
\[
B=(-5,5,-3)
\]
and
\[
C=(7,-1,9)
\]

Therefore the coordinates are $B=(-5,5,-3)$ and $C=(7,-1,9)$.

\item The plane through $A$, $B$ and $C$ has normal vector
\[
\overrightarrow{AB}\times\overrightarrow{AC}
\]

First find the two vectors:
\[
\overrightarrow{AB}=B-A=(-5-9,\,5-10,\,-3-(-1))=(-14,-5,-2)
\]
\[
\overrightarrow{AC}=C-A=(7-9,\,-1-10,\,9-(-1))=(-2,-11,10)
\]

Now
\begin{align*}
\overrightarrow{AB}\times\overrightarrow{AC}
&=
\begin{vmatrix}
\mathbf i & \mathbf j & \mathbf k\\
-14 & -5 & -2\\
-2 & -11 & 10
\end{vmatrix}\\
&=(-72,144,144)
\end{align*}
So a normal vector is
\[
(-1,2,2)
\]

Using the point $A=(9,10,-1)$, the plane equation is
\[
-(x-9)+2(y-10)+2(z+1)=0
\]
Simplifying,
\begin{align*}
-x+9+2y-20+2z+2&=0\\
-x+2y+2z-9&=0
\end{align*}
Hence
\[
x-2y-2z+9=0
\]

Therefore a Cartesian equation of the plane is $x-2y-2z+9=0$.

\item From part (a),
\[
\overrightarrow{AB}=(-14,-5,-2),
\qquad
\overrightarrow{AC}=(-2,-11,10)
\]
and from part (b),
\[
\overrightarrow{AB}\times\overrightarrow{AC}=(-72,144,144)
\]

Also
\[
D=(8,4,\alpha)
\]
so
\[
\overrightarrow{AD}=D-A=(8-9,\,4-10,\,\alpha-(-1))=(-1,-6,\alpha+1)
\]

The volume of tetrahedron $ABCD$ is
\[
V=\frac16\left|(\overrightarrow{AB}\times\overrightarrow{AC})\cdot\overrightarrow{AD}\right|
\]

So
\begin{align*}
(\overrightarrow{AB}\times\overrightarrow{AC})\cdot\overrightarrow{AD}
&=(-72,144,144)\cdot(-1,-6,\alpha+1)\\
&=72-864+144(\alpha+1)\\
&=72-864+144\alpha+144\\
&=144\alpha-648
\end{align*}

Given that the volume is $156$,
\[
\frac16|144\alpha-648|=156
\]
So
\[
|144\alpha-648|=936
\]
Divide by $72$:
\[
|2\alpha-9|=13
\]
Hence
\[
2\alpha-9=13
\quad\text{or}\quad
2\alpha-9=-13
\]
giving
\[
\alpha=11
\quad\text{or}\quad
\alpha=-2
\]

Therefore the possible values of $\alpha$ are $11$ and $-2$.

\item Since $\alpha>0$, from part (c) we use
\[
\alpha=11
\]
So
\[
D=(8,4,11)
\]

A direction vector for $l$ is obtained from its parametric form:
\[
(x,y,z)=\left(t-2,\frac{7-t}{2},t\right)
\]
so a direction vector is
\[
\left(1,-\frac12,1\right)
\]
and hence we can use
\[
\mathbf d_1=(2,-1,2)
\]

The line through $A$ and $D$ has direction vector
\[
\mathbf d_2=\overrightarrow{AD}=(8-9,4-10,11-(-1))=(-1,-6,12)
\]

Using $B=(-5,5,-3)$ on $l$, we have
\[
\overrightarrow{BA}=A-B=(14,5,2)
\]

For two skew lines,
\[
\text{shortest distance}=
\frac{\left|\,\overrightarrow{BA}\cdot(\mathbf d_1\times \mathbf d_2)\right|}{|\mathbf d_1\times \mathbf d_2|}
\]

Now find the cross product:
\begin{align*}
\mathbf d_1\times \mathbf d_2
&=
\begin{vmatrix}
\mathbf i & \mathbf j & \mathbf k\\
2 & -1 & 2\\
-1 & -6 & 12
\end{vmatrix}\\
&=(0,-26,-13)
\end{align*}
Its magnitude is
\[
|\mathbf d_1\times \mathbf d_2|
=\sqrt{0^2+(-26)^2+(-13)^2}
=\sqrt{845}
=13\sqrt5
\]

Also
\begin{align*}
\overrightarrow{BA}\cdot(\mathbf d_1\times \mathbf d_2)
&=(14,5,2)\cdot(0,-26,-13)\\
&=0-130-26\\
&=-156
\end{align*}
So the distance is
\begin{align*}
\frac{|-156|}{13\sqrt5}
&=\frac{156}{13\sqrt5}\\
&=\frac{12}{\sqrt5}\\
&=\frac{12\sqrt5}{5}
\end{align*}
and numerically
\[
\frac{12\sqrt5}{5}\approx 5.37
\]

Therefore the shortest distance is
\[
\frac{12\sqrt5}{5}
\]
which is $5.4$ to $2$ significant figures.
\end{enumerate}
\end{tcolorbox}


\newpage

% ---- Question 10 ----
\questionitem
% [Source: _QBT___Solns__Vector_Product, Q10]

The lines $L_1$, $L_2$ and $L_3$ are given by
\begin{align*}
L_1: &\quad \mathbf{r}=\begin{pmatrix}8\\3\\1\end{pmatrix}+\lambda\begin{pmatrix}3\\2\\-1\end{pmatrix}\\[6pt]
L_2: &\quad x=9-4\mu,\ y=-2+3\mu,\ z=-1+3\mu\\[6pt]
L_3: &\quad \dfrac{x-3}{1}=\dfrac{y+6}{-5}=\dfrac{z-1}{-2}
\end{align*}
The pairwise intersections of these lines are the vertices of a triangle. \\

Find the area of this triangle. \marks{9}

\vspace*{10pt}

\begin{tcolorbox}[colback=white, colframe=scarlet, title={\textbf{Solution}}, breakable, grow to left by=6mm, grow to right by=10mm]
Write the lines in parametric form:
\[
L_1:\ (x,y,z)=(8+3\lambda,\ 3+2\lambda,\ 1-\lambda)
\]
\[
L_2:\ (x,y,z)=(9-4\mu,\ -2+3\mu,\ -1+3\mu)
\]
and for \(L_3\), let the common value be \(t\). Then
\[
\frac{x-3}{1}=\frac{y+6}{-5}=\frac{z-1}{-2}=t
\]
so
\[
L_3:\ (x,y,z)=(3+t,\ -6-5t,\ 1-2t)
\]

Let the three vertices of the triangle be
\[
A=L_1\cap L_2,\qquad B=L_1\cap L_3,\qquad C=L_2\cap L_3
\]

For \(A=L_1\cap L_2\), equate coordinates:
\begin{align*}
8+3\lambda&=9-4\mu \\
3+2\lambda&=-2+3\mu \\
1-\lambda&=-1+3\mu
\end{align*}
From the third equation,
\[
2-\lambda=3\mu
\]
so
\[
\mu=\frac{2-\lambda}{3}
\]
Substitute into the first equation:
\begin{align*}
8+3\lambda&=9-\frac{4(2-\lambda)}{3} \\
24+9\lambda&=27-8+4\lambda \\
24+9\lambda&=19+4\lambda \\
5\lambda&=-5 \\
\lambda&=-1
\end{align*}
Hence
\[
\mu=\frac{2-(-1)}{3}=1
\]
So
\[
A=(8+3(-1),\ 3+2(-1),\ 1-(-1))=(5,1,2)
\]

For \(B=L_1\cap L_3\), equate coordinates:
\begin{align*}
8+3\lambda&=3+t \\
3+2\lambda&=-6-5t \\
1-\lambda&=1-2t
\end{align*}
From the third equation,
\[
\lambda=2t
\]
Substitute into the first equation:
\begin{align*}
8+3(2t)&=3+t \\
8+6t&=3+t \\
5t&=-5 \\
t&=-1
\end{align*}
Hence
\[
\lambda=2(-1)=-2
\]
So
\[
B=(8+3(-2),\ 3+2(-2),\ 1-(-2))=(2,-1,3)
\]

For \(C=L_2\cap L_3\), equate coordinates:
\begin{align*}
9-4\mu&=3+t \\
-2+3\mu&=-6-5t \\
-1+3\mu&=1-2t
\end{align*}
From the first equation,
\[
t=6-4\mu
\]
Substitute into the third equation:
\begin{align*}
-1+3\mu&=1-2(6-4\mu) \\
-1+3\mu&=1-12+8\mu \\
-1+3\mu&=-11+8\mu \\
10&=5\mu \\
\mu&=2
\end{align*}
Hence
\[
t=6-4(2)=-2
\]
So
\[
C=(9-4(2),\ -2+3(2),\ -1+3(2))=(1,4,5)
\]

Now use the formula for the area of a triangle:
\[
\text{Area}=\frac12\left|\mathbf{u}\times\mathbf{v}\right|
\]
where \(\mathbf{u}\) and \(\mathbf{v}\) are two side vectors.

Take
\[
\overrightarrow{AB}=
\begin{pmatrix}
2-5\\
-1-1\\
3-2
\end{pmatrix}
=
\begin{pmatrix}
-3\\
-2\\
1
\end{pmatrix},
\qquad
\overrightarrow{AC}=
\begin{pmatrix}
1-5\\
4-1\\
5-2
\end{pmatrix}
=
\begin{pmatrix}
-4\\
3\\
3
\end{pmatrix}
\]

Then
\[
\overrightarrow{AB}\times\overrightarrow{AC}
=
\begin{pmatrix}
(-2)(3)-(1)(3)\\
(1)(-4)-(-3)(3)\\
(-3)(3)-(-2)(-4)
\end{pmatrix}
=
\begin{pmatrix}
-9\\
5\\
-17
\end{pmatrix}
\]

So
\[
\left|\overrightarrow{AB}\times\overrightarrow{AC}\right|
=\sqrt{(-9)^2+5^2+(-17)^2}
=\sqrt{81+25+289}
=\sqrt{395}
\]

Therefore
\[
\text{Area of triangle}=\frac12\sqrt{395}
\]

So the area is \(\dfrac{\sqrt{395}}{2}\) square units, which is approximately \(9.94\) square units.
\end{tcolorbox}


\end{enumerate}
\end{document}
